\chapter{Derivative Validation}
\label{ch:bf-derivative-validation}

Derivative validation is a ladder.  Passing one rung is useful evidence, but it
does not certify the whole industrial target.

\section{Validation Ladder}
\label{sec:bf-derivative-validation-ladder}

BayesFilter derivative providers should pass:
\begin{enumerate}[label=(\roman*)]
  \item shape and finite-value checks;
  \item finite-difference checks on scalar or tiny LGSSMs;
  \item autodiff parity on small models where tape Hessians are feasible;
  \item backend parity among covariance, solve, Cholesky, and QR variants;
  \item factor reconstruction and solve-residual diagnostics;
  \item compiled/eager parity for TensorFlow backends;
  \item stress tests for low noise, collinearity, near rank loss, and masks.
\end{enumerate}

MacroFinance currently supplies examples of these tests, including QR
square-root parity against autodiff and finite differences, solve backend
agreement with the existing differentiated backend, TensorFlow eager/compiled
parity, and large-scale backend ladder diagnostics.  BayesFilter should import
the testing discipline before importing any claim of production readiness.

\section{Audit Record}
\label{sec:bf-derivative-audit-record}

Each derivative chapter should record the source labels, code files, and tests
that support it.  If a tool audit is inconclusive, the monograph should say so.
An inconclusive symbolic audit can still be valuable: it identifies which
guards, shapes, or noncommutative matrix steps need manual proof before a
formula is promoted to peer-review-ready status.
